\subsubsection{Autenticação e Autorização (RBAC)}
\label{sec:autenticacao}

A segurança de acesso é definida no \textit{backend} e refletida no cliente. No servidor, o fluxo de autenticação JWT compreende: (i) registro, com a senha \textit{hasheada} com Bcrypt (10 \textit{salt rounds}) antes da persistência; (ii) login, que valida credenciais e emite um \textit{access token} de curta duração (1 hora) e um \textit{refresh token} de longa duração (7 dias); (iii) refresh, em que o cliente troca o \textit{refresh token} por um novo par; e (iv) proteção de rotas, na qual o \texttt{JwtAuthGuard} valida o \textit{access token} e responde 401 quando inválido ou expirado.

O controle de acesso é composto por \textit{guards} aplicados conforme a natureza de cada rota, conforme a Tabela~\ref{tab:auth_pipeline}. O \texttt{JwtAuthGuard} é a base das rotas protegidas, pois constrói o \texttt{TenantContext} usado pelas demais validações. Os demais \textit{guards} são combinados nos controladores de acordo com a necessidade de papel, isolamento por organização ou acesso de desenvolvedor.

\begin{table}[H]
\centering
\caption{Guards de autorização do \textit{backend}}
\label{tab:auth_pipeline}
\begin{tabularx}{\textwidth}{l X}
\toprule
\textbf{Guard} & \textbf{Responsabilidade} \\
\midrule
\texttt{JwtAuthGuard}      & Valida o \textit{token} JWT e constrói o \texttt{TenantContext} em \texttt{req.context}. \\
\texttt{RolesGuard}        & Verifica o papel exigido por \texttt{@Roles()}; usuários \texttt{isDev} são dispensados. \\
\texttt{TenantFilterGuard} & Confere se o parâmetro \texttt{:organizationId} da rota coincide com o do \texttt{req.context}. \\
\texttt{DevGuard}          & Restringe rotas \texttt{@Dev()} a usuários com \texttt{isDev=true}. \\
\midrule
\multicolumn{2}{l}{\textit{Aprovado nos guards aplicáveis} $\rightarrow$ a requisição alcança o Controller.} \\
\bottomrule
\end{tabularx}
\fonte{Elaborado pelo autor (2026)}
\end{table}

O \texttt{JwtAuthGuard} enriquece o \textit{payload} JWT no momento do login com todos os dados necessários (\texttt{userId}, \texttt{organizationId}, \texttt{role}, \texttt{isDev}), eliminando consultas ao banco a cada requisição. Usuários listados em \texttt{DEV\_EMAILS} recebem \texttt{isDev=true} no token e ignoram validações de tenant e role. Em termos de papéis, o ADMIN tem acesso total (organizações, usuários, motoristas, configurações), o DRIVER tem acesso limitado (visualizar viagens e confirmar ações relacionadas) e o USER implícito (passageiro) pode inscrever-se em viagens e visualizar suas inscrições. As \textit{memberships} permitem múltiplas \textit{roles} por usuário em diferentes organizações (ex.: motorista em uma org, admin em outra).

Reflexo no cliente. A aplicação persiste a sessão em \texttt{localStorage} (\textit{token} de acesso, \textit{token} de renovação e dados básicos do usuário) e provê o estado de autenticação por meio de dois contextos encadeados: um de autenticação, que envolve um de papel, garantindo que a determinação de papel só ocorra quando a identidade já está disponível. A principal característica de resiliência é a renovação transparente de sessão, concentrada no cliente HTTP único (\texttt{lib/api.ts}): por padrão, cada requisição anexa o cabeçalho \texttt{Authorization: Bearer}; diante de uma resposta \texttt{401}, o cliente dispara uma única renovação do \textit{access token}, deduplicada entre requisições concorrentes por meio de uma \textit{promise} compartilhada, e, em caso de sucesso, repete automaticamente a chamada original uma única vez. Se a renovação falhar, o erro original é propagado e a sessão é encerrada. Toda falha é, ainda, normalizada em uma exceção tipada (\texttt{ApiError}) que carrega \texttt{status}, dados e \texttt{errorCode}.
